% © 2025 Ing. David Jaroš — CC BY-NC-ND 4.0
%
% This work is licensed under a Creative Commons Attribution-NonCommercial-NoDerivatives
% 4.0 International License (CC BY-NC-ND 4.0).
%
% License History: Earlier drafts (up to v0.3) were released under CC BY 4.0.
% From v0.4 onward, all material is released under CC BY-NC-ND 4.0 to protect
% the integrity of the theoretical work during ongoing academic development.
%
% See LICENSE.md for full license text.

\documentclass[11pt,a4paper]{article}
\usepackage{amsmath, amssymb, amsthm}
\usepackage{hyperref}
\usepackage[a4paper, margin=2.5cm]{geometry}
\usepackage{tcolorbox}
\usepackage{longtable}
\usepackage{booktabs}

\newtheorem{theorem}{Theorem}[section]
\newtheorem{proposition}[theorem]{Proposition}
\newtheorem{lemma}[theorem]{Lemma}
\newtheorem{corollary}[theorem]{Corollary}
\theoremstyle{definition}
\newtheorem{definition}[theorem]{Definition}
\theoremstyle{remark}
\newtheorem{remark}[theorem]{Remark}

\title{\textbf{Theta Spectral Framework}\\
\large Unified Biquaternion Theory — Canonical Appendices Series}
\author{Ing.\ David Jaroš}
\date{March 2026}

\begin{document}
\maketitle
\tableofcontents
\bigskip

\begin{tcolorbox}[colback=blue!4!white,colframe=blue!50!black,title=Purpose]
This appendix formalises the theta spectral representation of the
$\Theta$-field used across UBT.  It defines the spectral decomposition,
connects toroidal topology with Jacobi theta functions, and analyses the
role of the modular parameter $\tau = t + i\psi$.

\medskip
\textbf{Primary sources}:
\texttt{canonical/UBT\_canonical\_main.tex \S10},
\texttt{canonical/geometry/Rpsi\_dynamical\_fix.tex},
\texttt{research\_tracks/research/partition\_function\_modular.md}

\medskip
\begin{tabular}{lll}
\hline
Item & Status & Section \\
\hline
$\Theta$ spectral decomposition & \textbf{Defined} & \S\ref{sec:spectral_decomp} \\
Toroidal topology $T^3\times S^1_\psi$ & \textbf{Proved} & \S\ref{sec:torus} \\
Jacobi theta functions & \textbf{Connected} & \S\ref{sec:jacobi} \\
Modular parameter $\tau$ & \textbf{Analysed} & \S\ref{sec:modular} \\
Poisson summation duality & \textbf{Proved} & \S\ref{sec:poisson} \\
\hline
\end{tabular}
\end{tcolorbox}

% -----------------------------------------------------------------
\section{Theta Spectral Decomposition}\label{sec:spectral_decomp}
% -----------------------------------------------------------------

\begin{definition}[Spectral decomposition of $\Theta$]
\label{def:spectral}
On the torus $T^3 \times S^1_\psi$ with modular parameter $\tau = t + i\psi$,
the biquaternion field $\Theta(q,\tau)$ admits the spectral decomposition:
\begin{equation}\label{eq:spectral}
  \Theta(q,\tau)
  = \sum_{n \in \mathbb{Z}}\sum_{\mathbf{k}\in\mathbb{Z}^3}
    \hat{\Theta}_{n,\mathbf{k}}\;
    e^{i n\psi/R_\psi}\,
    e^{i \mathbf{k}\cdot\mathbf{x}/R_t},
\end{equation}
where $n$ is the winding number on $S^1_\psi$, $\mathbf{k} = (k_1,k_2,k_3)$
are the spatial Kaluza--Klein momenta, $R_\psi$ and $R_t$ are the respective
radii, and $\hat{\Theta}_{n,\mathbf{k}} \in \mathbb{C}\otimes\mathbb{H}$
are the spectral coefficients.
\end{definition}

\begin{remark}
In the isotropic limit $R_t \to \infty$ (decompactification of spatial dimensions),
the spatial sum reduces to a continuous Fourier transform, recovering the
standard momentum space representation.
\end{remark}

\subsection{Mode Classification}

\begin{center}
\begin{tabular}{cll}
\toprule
$n$ & Mode type & Physical interpretation \\
\midrule
$n = 0$ & Zero mode & Vacuum / background field \\
$n = 1$ & Fundamental winding & Ground-state charged excitation \\
$|n| > 1$ & Higher winding & Higher-charge or heavier excitations \\
$n = n^* = 137$ & Prime attractor & Fine-structure mode (minimum of $V_{\mathrm{eff}}$) \\
\bottomrule
\end{tabular}
\end{center}

% -----------------------------------------------------------------
\section{Toroidal Topology}\label{sec:torus}
% -----------------------------------------------------------------

\begin{proposition}[Geometry of compactification]\label{prop:torus}
The UBT field lives on the total space
$\mathcal{M} \times T^3 \times S^1_\psi$, where:
\begin{itemize}
  \item $\mathcal{M} = \mathbb{R}^{1,3}$ is the non-compact Minkowski spacetime
    (or its curved generalisation)
  \item $T^3 = (S^1_{R_t})^3$ is the compactified spatial torus
  \item $S^1_\psi = S^1_{R_\psi}$ is the imaginary-time circle with $\psi \sim \psi + 2\pi R_\psi$
\end{itemize}
\end{proposition}

The partition function of the $T^3$ sector is:
\begin{equation}\label{eq:Z_T3}
  \hat{Z}_{T^3}(\tau)
  = \sum_{\mathbf{k}\in\mathbb{Z}^3} q^{|\mathbf{k}|^2},
  \qquad q = e^{2\pi i\tau},
\end{equation}
which equals the cube of the Jacobi theta function:
\begin{equation}\label{eq:Z_theta3}
  \hat{Z}_{T^3}(\tau) = \vartheta_3(\tau)^3.
\end{equation}

\noindent\textbf{Source}: \texttt{research\_tracks/research/partition\_function\_modular.md}.

% -----------------------------------------------------------------
\section{Jacobi Theta Functions and Toroidal Topology}\label{sec:jacobi}
% -----------------------------------------------------------------

\begin{definition}[Jacobi theta function $\vartheta_3$]
\begin{equation}\label{eq:theta3}
  \vartheta_3(\tau)
  = \sum_{n\in\mathbb{Z}} e^{i\pi n^2\tau}
  = 1 + 2\sum_{n=1}^\infty e^{i\pi n^2\tau},
  \qquad \mathrm{Im}(\tau) > 0.
\end{equation}
\end{definition}

\begin{theorem}[Connection to UBT partition function]\label{thm:partition}
The partition function of the compactified $\Theta$ field on $T^3$
is $\hat{Z}_{T^3}(\tau) = \vartheta_3(\tau)^3$, transforming under
modular inversion as:
\begin{equation}\label{eq:modular_transform}
  \hat{Z}_{T^3}(-1/\tau)
  = (-i\tau)^{3/2}\,\hat{Z}_{T^3}(\tau).
\end{equation}
The modular weight is $k = 3/2 = d/2$ where $d = 3 = \dim_\mathbb{R}(\mathrm{Im}\,\mathbb{H})$.
\end{theorem}

\begin{proof}
The Jacobi identity for $\vartheta_3$ gives:
\begin{equation}
  \vartheta_3(-1/\tau) = (-i\tau)^{1/2}\,\vartheta_3(\tau).
\end{equation}
Cubing: $\vartheta_3(-1/\tau)^3 = (-i\tau)^{3/2}\,\vartheta_3(\tau)^3$.
The weight $3/2 = d/2$ follows from $d = 3 = \dim_\mathbb{R}(\mathrm{Im}\,\mathbb{H})$,
giving the same algebraic origin as the $B_{\mathrm{base}} = N_{\mathrm{eff}}^{3/2}$
exponent in the $\alpha$ derivation. $\square$
\end{proof}

\noindent\textbf{Source}: \texttt{research\_tracks/research/partition\_function\_modular.md \S v57--v59}.

\begin{remark}
The modular weight $k = 3/2$ of $\hat{Z}_{T^3}(\tau)$ is consistent with the
$B_{\mathrm{base}} = N_{\mathrm{eff}}^{3/2}$ exponent in the $\alpha$-derivation.
Both arise from $d = \dim_\mathbb{R}(\mathrm{Im}\,\mathbb{H}) = 3$.
However, the Kac-Moody level $k_{\mathrm{KM}} = 1$ (Gap G3-k) is a \emph{different}
quantity and remains open.
\end{remark}

% -----------------------------------------------------------------
\section{Role of the Modular Parameter $\tau$}\label{sec:modular}
% -----------------------------------------------------------------

\subsection{Complex Time as Modular Parameter}

In UBT, the complex time $\tau = t + i\psi$ plays the dual role of:
\begin{enumerate}
  \item \textbf{Physical time parameter}: $t$ is real Minkowski time; $\psi$
    is the imaginary part controlling quantum phase evolution.
  \item \textbf{Modular parameter} of the toroidal compactification: $\tau$
    parametrises the shape (complex structure) of the torus $T^2$ formed
    by the time circle and the imaginary-time circle.
\end{enumerate}

\begin{tcolorbox}[colback=yellow!5!white,colframe=orange!60!black,title=Important: $\psi$ is NOT Wick rotation]
$\psi \in \mathbb{R}$ is a \emph{dynamical} degree of freedom — the imaginary
component of complex time.  It is \emph{not} a Wick rotation (which is a formal
analytic continuation).  This is proved in
\texttt{canonical/THEORY/math/fields/biquaternion\_time.tex \S5.1}.
\end{tcolorbox}

\subsection{Modular Group Action}

The modular group $\mathrm{SL}(2,\mathbb{Z})$ acts on $\tau$ by
\begin{equation}
  \tau \;\mapsto\; \frac{a\tau + b}{c\tau + d},
  \qquad
  \begin{pmatrix}a&b\\c&d\end{pmatrix} \in \mathrm{SL}(2,\mathbb{Z}).
\end{equation}
The UBT partition function $\hat{Z}_{T^3}(\tau) = \vartheta_3(\tau)^3$ transforms
as a modular form of weight $3/2$ under the index-2 subgroup of
$\mathrm{SL}(2,\mathbb{Z})$ fixing $\vartheta_3$.

\subsection{Spectral Modes and Hecke Operators}

The Hecke operator $T_p$ acts on the spectral coefficients $\hat{\Theta}_{n,\mathbf{k}}$
as:
\begin{equation}
  (T_p\,\hat{Z})_n
  = \sum_{d|n,\; d|p} d^{k-1}\hat{Z}_{n^2/(dp)},
\end{equation}
and the prime $p = 137$ is a Hecke eigenvalue associated with the attractor mode
$n^* = 137$ (numerical support: \texttt{docs/reports/hecke\_lepton/}).

% -----------------------------------------------------------------
\section{Poisson Summation Duality}\label{sec:poisson}
% -----------------------------------------------------------------

\begin{theorem}[Poisson summation duality]\label{thm:poisson}
For the Gaussian function $f(\mathbf{k}) = e^{-\pi|\mathbf{k}|^2/\tau}$
on $\mathbb{Z}^d$, the winding (momentum) sum and the spectral (winding) sum
are related by:
\begin{equation}\label{eq:poisson}
  W(\tau)
  := \sum_{\mathbf{k}\in\mathbb{Z}^d} e^{-\pi|\mathbf{k}|^2/\tau}
  = \tau^{d/2}\,S(\tau),
\end{equation}
where $S(\tau) = \sum_{\mathbf{n}\in\mathbb{Z}^d} e^{-\pi|\mathbf{n}|^2\tau}$
is the spectral sum.  For $d=3$, the prefactor is $\tau^{3/2}$.
\end{theorem}

\begin{proof}
Standard Poisson summation formula:
$\hat{f}(\mathbf{n}) = \tau^{d/2}e^{-\pi\tau|\mathbf{n}|^2}$
for $f(\mathbf{k}) = e^{-\pi|\mathbf{k}|^2/\tau}$.
Summing over $\mathbb{Z}^d$ gives~\eqref{eq:poisson}. $\square$
\end{proof}

\noindent\textbf{Source}: \texttt{ubt\_with\_chronofactor/scripts/spectral/poisson\_duality\_demo.py}.

\begin{remark}
The prefactor $\tau^{d/2}$ with $d = 3$ is the same $3/2$ exponent as the
modular weight of $\hat{Z}_{T^3}$ and the $B_{\mathrm{base}}$ exponent.
All three arise from $d = \dim_\mathbb{R}(\mathrm{Im}\,\mathbb{H}) = 3$.
\end{remark}

% -----------------------------------------------------------------
\section{Summary}\label{sec:summary_theta}
% -----------------------------------------------------------------

\begin{center}
\begin{longtable}{llll}
\toprule
Result & Status & Layer & Source \\
\midrule
Spectral decomposition~\eqref{eq:spectral} & Defined & L0 &
  Def.~\ref{def:spectral} \\
Toroidal topology $T^3\times S^1_\psi$ & Proved & L0 &
  Prop.~\ref{prop:torus} \\
$\hat{Z}_{T^3} = \vartheta_3(\tau)^3$ & Proved & L0 &
  Thm.~\ref{thm:partition} \\
Modular weight $k = 3/2 = d/2$ & Proved & L0 &
  Thm.~\ref{thm:partition} \\
Poisson duality $W = \tau^{3/2}S$ & Proved & L0 &
  Thm.~\ref{thm:poisson} \\
$\psi$ dynamical (not Wick rotation) & Proved & L0 &
  \S\ref{sec:modular} \\
Hecke eigenvalue $p=137$ & Numerical support & L1 &
  \texttt{docs/reports/hecke\_lepton/} \\
Kac-Moody level $k_{\mathrm{KM}}=1$ (Gap G3-k) & Open & L1--L2 &
  \texttt{DERIVATION\_INDEX.md} \\
\bottomrule
\end{longtable}
\end{center}

\end{document}
